\begin{claim}{D.71}{Proved}{The Frobenius cover of a scalar cubic torus coordinate}
Let $a$ be an order-three element of the full
once-punctured-torus loop-outer group, with actual
order-three representative $\alpha$ and labelled
extension $E=F\rtimes_\alpha\langle s\rangle$.
At every prime $q\ne3$ with
$\lambda_q\ne1$, there is the following actual
finite quotient, after the explicit based peripheral
change of generators in the proof:
\[
 \begin{gathered}
 Q=\langle r,t\mid r^q=t^3=1,\ trt^{-1}=r^\ell\rangle,
 \qquad
 \ell=\lambda_q\bmod q\in\F_q^\times,\\
 \rho:F\twoheadrightarrow Q,\qquad
 \rho(x)=t,\quad\rho(y)=r,\quad
 \rho\alpha=\rho .
 \end{gathered}                                                \tag{D.71a}
\]
Here $\ell$ has order three and $|Q|=3q$.
Put $R=\langle r\rangle$, $N=\ker\rho$,
$C=C_F(s)$, and $C_0=C\cap N$.
The actual cover defining $N$ has type
$((3q-1)/2,3)$, and
\[
 1\longrightarrow C_0\longrightarrow C
      \xrightarrow{\ \rho|_C\ }R\longrightarrow1,
 \qquad
                       H^1(C_0,\F_3)=0 .                    \tag{D.71b}
\]
There are exactly three $N$-conjugacy classes of
order-three sections in $N\rtimes\langle s\rangle$.
With their original quotient-generator labels,
their representatives and exact normalizers are
\[
 \begin{gathered}
 T_i=x^i\langle s\rangle x^{-i},\\
 N_{N\rtimes\langle s\rangle}(T_i)
       =x^i(C_0\times\langle s\rangle)x^{-i}
                      \qquad(0\leq i<3).
 \end{gathered}                                                \tag{D.71c}
\]
On the integral character module
$V_3=H^1(N,\Z_3)$, retain the actual rotation
projector and its complement
\[
 e=\frac1q\sum_{g\in R}g,\qquad W_3=(1-e)V_3 .
                                                               \tag{D.71d}
\]
For the action $S$ given by pullback under
$\alpha^{-1}$, the module $W_3$ is free of rank
$q-1$ over $\Z_3[\langle S\rangle]\cong\Z_3[C_3]$.
The isomorphism is supplied by lifting cyclic
vectors for the actual action.
The averaged rank-four module $eV_3$ reduces to
two Jordan blocks of size two modulo three.
These are exact quotient, lattice, and fixed-group
constraints. No exclusion or realization of the
scalar cubic coordinate is asserted.
\end{claim}
\begin{proof}
\emph{The actual index-three cover and its based generators.}
Use D.48 and D.51 for the two section classes of
$E$, with $\langle s\rangle$ the chosen one.
There is an invariant character $\chi:F\to\F_3$,
extended by $\chi(s)=0$, which can be scaled
uniquely so that
\[
 \chi(v)=1,\qquad
 \chi(s_1)=-1,\qquad
 vcv^{-1}=c^\lambda,\qquad v^3=c^\kappa ,
 \quad c=[X,Y],
                                                               \tag{D.71e}
\]
where $v$ is a cusp-normalizer element labelled by
$s$, and $s_1$ is the labelled generator of the
other full section class.
Indeed D.51 gives opposite nonzero values for
the cusp line and the other section line.
Also $\kappa_3$ is a unit by D.48.

Here is an explicit peripheral change of basis
with $\chi(x)=1,\chi(y)=0$, without rescaling
the now specified character.
For any current basis $(X,Y)$ the transformation
\[
 Q_0(X,Y)=(XYX^{-1},X^{-1}),\qquad
 Q_0^{-1}(x,y)=(y^{-1},yxy^{-1})
                                                               \tag{D.71f}
\]
has the displayed inverse and preserves
$[X,Y]$ exactly. Its action on the two character
values is $(a,b)\mapsto(b,-a)$.
Among its four iterates, at least one has first
character value $1$, since $(a,b)\ne(0,0)$.
Choose the first such iterate $(x,Y')$,
and let $d\in\{0,1,-1\}$ represent $\chi(Y')$.
Set $y=Y'x^{-d}$.
This shear has inverse $Y'=yx^d$ and
$[x,Y'x^{-d}]=[x,Y']$.
These actual discrete free-group automorphisms,
with the stated inverses, therefore supply
the required basis in the full completion
and keep $c=[x,y]$ exactly.

Let $F_0=\ker\chi$.
Its actual type-$(1,3)$ covering has the
Schreier basis and peripheral representatives
\[
 \begin{gathered}
 u=x^3,\qquad B_i=x^iyx^{-i}\quad(0\leq i<3),\\
 c_i=x^icx^{-i}=B_{i+1}B_i^{-1},\qquad
 B_3=uB_0u^{-1}.
 \end{gathered}                                                \tag{D.71g}
\]
These identities are actual words, including
the final value $i=2$.
On $M=H_1(F_0,\widehat\Z)$, deck conjugation by
$x$ gives $J$, which fixes $u$ and cyclically
permutes the $B_i$.
The inclusion homology map is the actual map
\[
 I:M\longrightarrow H_1(F,\widehat\Z),\qquad
 I(u)=3[x],\quad I(B_i)=[y].
                                                               \tag{D.71h}
\]
Its kernel is the rank-two boundary augmentation,
generated by the three differences $B_{i+1}-B_i$
with their sole sum relation.
Filling the three punctures gives the actual
degree-three closed torus cover; (D.71h) factors
through its integral lattice inclusion.

Let $S_{\rm hom}$ be induced by $\alpha$ on $M$.
Since $\chi\alpha=\chi$, it commutes with $J$:
the conjugating elements $\alpha(x)$ and $x$
differ by an element of $F_0$, whose inner
action on $M$ is trivial.
Write $v=f_0s$. Equation (D.71e) gives
$\chi(f_0)=1$, so the homology action $T$ of
conjugation by $v$ is
\[
                              T=JS_{\rm hom}.                \tag{D.71i}
\]
It fixes each of the three peripheral classes
with multiplier $\lambda$.
To check its labels, conjugation by $v$ preserves
$\chi$ on $F$, so $vx^iv^{-1}$ differs from
$x^i$ by an element of $F_0$, while
$vcv^{-1}=c^\lambda$.
Thus
\[
       S_{\rm hom}|_{\ker I}=\lambda J^{-1}.
                                                               \tag{D.71j}
\]

There is also a fully specified matrix for this
comparison. If $A=(a_{ij})$ is the original
homology matrix in the basis $(x,y)$ and
$\boldsymbol1=(1,1,1)^{\mathsf T}$, then
\[
 S_{\rm hom}=
 \begin{pmatrix}
 a_{11}&(a_{12}/3)\boldsymbol1^{\mathsf T}\\
 a_{21}\boldsymbol1&
       \lambda J_3^{-1}
          +((a_{22}-\lambda)/3)\boldsymbol1\boldsymbol1^{\mathsf T}
 \end{pmatrix},
                                                               \tag{D.71k}
\]
where $J_3$ is the cycle on $(B_0,B_1,B_2)$.
The numerator divisibilities hold at $3$:
$\chi\alpha=\chi$ gives $a_{11}\equiv1$ and
$a_{12}\equiv0$ modulo three; the determinant
is $\lambda$, with $\lambda_3=1$, so
$a_{22}\equiv1\equiv\lambda$ there.
At every other prime division by three is
invertible. To prove the formula, the $u$-column
is determined by $J$-invariance and (D.71h);
the remaining three columns are determined by
(D.71j), commutation with $J$, and
$IS_{\rm hom}=AI$.
These determinations can first be read over each
$\Q_p$; the stated integral divisibilities
and torsion-free lattices give (D.71k) in
the full module.

\emph{The invariant character and the actual Frobenius quotient.}
Fix the specified bad prime $q$.
The unit $\lambda_q$ is a primitive cube root.
It reduces to one modulo $q$ only for the
identity cube root, because $q\ne3$ and the
three roots have distinct reductions.
Thus $\ell=\lambda_q\bmod q$ has order three
and $q\equiv1\pmod3$.
Define the continuous character
\[
 \ell_q:F_0\longrightarrow\Z_q,\qquad
           \ell_q(u)=0,\quad
           \ell_q(B_i)=\lambda_q^i .
                                                               \tag{D.71l}
\]
Its values sum to zero, and
$\ell_qJ=\lambda_q\ell_q$.
Over $\Z_q$, the row therefore annihilates the
deck-trivial summand and is supported on the
boundary augmentation. Equation (D.71j), or directly
(D.71k), gives the exact equation
\[
                             \ell_q\alpha=\ell_q.             \tag{D.71m}
\]
Put $N=\ker(\ell_q\bmod q)$ inside $F_0$.
It is normal in $F$: conjugation by $x$ scales
the character by the unit $\ell$, and
conjugation by $F_0$ does not change a character.
It is also $\alpha$-invariant in both directions
by (D.71m).
Its quotient in $F_0$ has order $q$, so
$[F:N]=3q$.

In $F/N$, the images of $x,y$ satisfy
$x^3=u=1$, $y^q=1$, and
$xyx^{-1}=y^\ell$.
The latter equality is the equality of their
images under $\ell_q\bmod q$ in $F_0/N$.
The presented group in (D.71a) has the explicit
multiplication
\[
 (a,i)(b,j)=(a+\ell^ib,\ i+j)
         \quad\text{on }\F_q\times\F_3,
 \qquad r=(1,0),\quad t=(0,1).
                                                               \tag{D.71n}
\]
It has order $3q$.
The inverse coordinate homomorphisms are
$r\mapsto yN,t\mapsto xN$ and
$xN\mapsto t,yN\mapsto r$.
The relations and the already computed index
make them inverse, proving exactly (D.71a)'s
group and quotient, rather than just a
surjection to a smaller group.

The induced automorphism on $Q$ fixes $r$
by (D.71m) and preserves the quotient label
in $Q/R$ by $\chi\alpha=\chi$.
It consequently sends $t$ to $r^jt$ for
some $j\in\F_q$.
Its third iterate sends $t$ to $r^{3j}t$.
Since the original $\alpha$ has order three
and $q\ne3$, $j=0$.
Thus $\rho\alpha=\rho$.
The actual extension quotient and inverse are
\[
 \begin{aligned}
 E/N&\longrightarrow Q\times C_3,
               &fs^iN&\longmapsto(\rho(f),s^i),\\
 Q\times C_3&\longrightarrow E/N,
               &(z,s^i)&\longmapsto fs^iN
                         \quad(\rho(f)=z).
 \end{aligned}                                                \tag{D.71o}
\]
Independence of the lift and multiplication
follow from $\rho\alpha=\rho$.

The cusp has image
$\rho(c)=[t,r]=r^{\ell-1}$ of order $q$.
Its action on the $3q$ cosets has three cycles
of length $q$.
The corresponding discrete cover therefore
has three punctures and Euler characteristic
$-3q$, hence genus $(3q-1)/2$.
Its finite-index subgroup carries its full
induced profinite topology, giving exactly
the stated full surface completion $N$.

\emph{The integral rotation module, including its inverse maps.}
Write $\Lambda_3=\Z_3[R]$ and
$\nu=\sum_{i=0}^{q-1}r^i$.
Let $a,b\in\{1,\ldots,q-1\}$ be the integer
representatives of $\ell,\ell^2$, and put
$P_k=1+r+\cdots+r^{k-1}$.
The actual degree-$q$ graph cover of the
four-generator rose for $F_0$ has boundary
\[
 \partial:\Lambda_3^4\longrightarrow\Lambda_3,\qquad
 (A,B,D,E)\longmapsto
       (r-1)B+(r^a-1)D+(r^b-1)E .
                                                               \tag{D.71p}
\]
Its exact cycle-coordinate maps are
\[
 \begin{aligned}
 \Lambda_3^3\oplus\Z_3&\longrightarrow\ker\partial,
 &(A,D,E,k)&\longmapsto
          (A,-P_aD-P_bE+k\nu,D,E),\\
 \ker\partial&\longrightarrow\Lambda_3^3\oplus\Z_3,
 &(A,B,D,E)&\longmapsto(A,D,E,k),\\[-2pt]
 &&&\hspace{-15mm}B+P_aD+P_bE=k\nu .
 \end{aligned}                                                \tag{D.71q}
\]
The scalar $k$ exists uniquely: the annihilator
of $r-1$ consists exactly of constant-coefficient
multiples of $\nu$, as comparison of the cyclic
coefficients proves.
Thus these are inverse integral module maps and
\[
                  H_1(N,\Z_3)\cong
                             \Z_3\oplus\Lambda_3^3 .
                                                               \tag{D.71r}
\]
Dualizing gives the same deck multiplicities
in $V_3$.
Restriction and transfer for the actual
index-$q$ subgroup $N\subset F_0$ identify
$eV_3$ with $H^1(F_0,\Z_3)$.
The maps are equivariant under the original
$\alpha$; the remaining module $W_3$
has rank $3(q-1)$.

On $H_1(F_0,\Q_3)$, the boundary part of
$S_{\rm hom}$ is $J^{-1}$ by (D.71j).
Its two eigenvalues are the primitive cube
roots. The filled rank-two quotient has
the original nonidentity cubic matrix
$A_3$ by the lattice map (D.71h).
D.48 proves its polynomial is $X^2+X+1$.
Consequently the same polynomial annihilates
the whole rank-four module:
first over $\Q_3$, by the two commuting
$J$-summands, and then over $\Z_3$ by
torsion-freeness.
For the dual convention
$S=(S_{\rm hom}^{-1})^{\mathsf T}$ this gives
\[
                        I+S+S^2=0\quad\text{on }eV_3.
                                                               \tag{D.71s}
\]
The ring $\Z_3[z]/(z^2+z+1)$ is a discrete
valuation ring: with $\pi=z-1$ its defining
polynomial is $\pi^2+3\pi+3$, and
$(3)=(\pi^2)$.
The rank-four $\Z_3$-lattice in (D.71s)
is a torsion-free module over this ring,
so is free of rank two.
Its reduction modulo three is therefore
exactly two Jordan blocks of size two.

\emph{The integral fixed-cup obstruction at three.}
We use the following calculation with the
actual pro-three closed-surface group.
Let $P$ be such a group of genus at least two
and let $b$ be an actual automorphism of
order three whose integral homology action
is nonidentity.
Any two $b$-fixed characters
$\theta,\psi:P\to\Z_3$ have cup pairing
divisible by three.

Indeed put $G=P\rtimes_b\langle\tau\rangle$.
Both characters extend by
$\widetilde\theta(h\tau^i)=\theta(h)$ and
$\widetilde\psi(h\tau^i)=\psi(h)$.
They kill every subgroup of order three,
since $\Z_3$ is torsion-free.
Such a subgroup $T=\langle n\tau\rangle$
has the exact normalizer
$C_P(T)\times T$, by splitting off its
unique element with a prescribed quotient
label, as in (D.70e).
Its actual action on $P$ is a nonidentity
inner twist of $b$, because inner twists
do not change the given nonidentity
integral homology matrix.
The Demushkin conclusion of
\cite[Corollary~2 and \S2]{A-KochloukovaZalesskii2007}
therefore makes $C_P(T)$ free pro-three.
The restricted product of the two character
classes is zero on every normalizer.

If their reduced cup were a nonzero surface
class $\eta$, the actual product of the
extensions would lift $\eta$ to $H^2(G,\F_3)$.
It is consequently a permanent nonzero
class in the highest row $E_2^{0,2}$
of the Hochschild--Serre spectral sequence.
The action on $H^2(P,\F_3)$ is identity,
since an order-three scalar in $\F_3^\times$
is identity.
Let $z\in H^2(C_3,\F_3)$ be its period
generator. Its inflation is permanent by
the actual section, so $z\eta$ is a
nonzero permanent element of $E_2^{2,2}$.
There is no incoming differential to
this highest row, and its outgoing
differentials vanish by multiplicativity.
The class
$z\overline{\widetilde\theta}
       \overline{\widetilde\psi}$ is therefore
nonzero in $H^4(G,\F_3)$.
Its restrictions to all normalizers are zero.
This contradicts Brown injectivity in
degree four
\cite[Theorem~2]{A-KochloukovaZalesskii2007}:
$P$ is open of cohomological dimension two,
and finite subgroups of $G$ inject into
$C_3$, so there is no $C_3\times C_3$.
This proves the divisibility assertion.
It requires no assumption that $b$ is
identity on mod-three homology.

\emph{The actual filled character module and its paired lattices.}
Fill the three actual cusps
$d_i=x^ic^qx^{-i}$ of $N$ and put
\[
 P=\left(N/
       \overline{\langle\!\langle d_0,d_1,d_2\rangle\!\rangle}
                   \right)^{(3)},\qquad
 p:N\twoheadrightarrow P .
                                                               \tag{D.71t}
\]
The same finite-target universal property
as in (D.70i) identifies the quotient
before maximal pro-three completion with
the full completion of the closed covering
surface. Thus $P$ is the genus-$(3q-1)/2$
closed-surface pro-three group.
The normal closure is $\alpha$-invariant
with the original labels: for
$g=x^i$, write
$\alpha(g)f_0^{-1}=n x^j c^k$
with $n\in N$, choosing $j\in\{0,1,2\}$
by its coset in $Q/R$ and $k$ by
the rotation coordinate.
This is possible because $\rho(c)$
generates $R$.
Then
$\alpha(gc^qg^{-1})=n d_j^\lambda n^{-1}$.
The inverse proves equality of the
closed normal subgroups.
Hence $\alpha$ induces an actual
automorphism of $P$.

Pullback gives the exact dual injection
$H^1(P,\Z_3)\hookrightarrow V_3$ whose
image is the characters vanishing on
all $d_i$: such a character factors
through the full quotient and then
through its maximal pro-three quotient,
and the converse follows from (D.71t).
The three cusp evaluations are
$R$-invariant, since the rotations
fix their labels and their positive
generators up to conjugation.
Thus $W_3$ lies in this closed character
module. Its complement there is the
averaged rank-two closed-torus module.
The original covering homology map identifies
the action on the filled rank-two homology
over $\Q_3$ with the nonidentity matrix $A_3$.
On its dual the pullback convention is
$(A_3^{-1})^{\mathsf T}$ after that rational
identification. It is nonidentity as well.
Thus the induced automorphism of $P$ has
order exactly three and is nonidentity
on integral homology.

The rotation action extends over the
three filled punctures and preserves
the integral surface orientation.
The odd projector $e$ is therefore
self-adjoint for the unimodular
alternating cup form, exactly as in
(D.70p), with $q$ a unit in $\Z_3$.
Consequently its complementary closed
summand $W_3$ has a unimodular
alternating form.
The automorphism $S$ preserves this
form: its action on the integral
surface class has order dividing three
in $\Z_3^\times$, which has no
nonidentity element of order three.

Choose a finite unramified extension
$\mathcal O/\Z_3$ containing the $q$th
roots of unity.
For a character $\psi:R\to\mathcal O^\times$,
the actual projector is
$e_\psi=q^{-1}\sum_{g\in R}\psi(g)^{-1}g$.
Each nontrivial character lattice
\[
                 V_\psi=e_\psi(W_3\otimes\mathcal O)
                                                               \tag{D.71u}
\]
has rank three by (D.71r).
The automorphism $S$ commutes with
every rotation operator because
$\rho\alpha=\rho$, so preserves
each such lattice.
The cup form pairs $V_\psi$
perfectly with $V_{\psi^{-1}}$
and pairs it trivially with every
other character lattice.
The inverse character is distinct,
since $q$ is odd.

Let $K_3$ be the fraction field of
$\mathcal O$.
The polynomial $X^2+X+1$ is irreducible
over $K_3$: a square root of $-3$
would have half-integral valuation
in this unramified extension.
Since $S^3=1$, its fixed subspace
on the rank-three $K_3$-space has
dimension one or three.
The fixed-cup divisibility just proved
holds after extension to $\mathcal O$.
Indeed $\mathcal O$ is finite free
over $\Z_3$, so kernels of $S-I$
commute with this flat scalar
extension, and the bilinear divisibility
passes to their scalar extensions.

Fixed dimension three is impossible.
It would make $S=I$ on $V_\psi$,
and preservation of the perfect paired
form would give $S=I$ on
$V_{\psi^{-1}}$.
The perfect integral pairing would
then give two fixed characters with
unit cup pairing, contradicting that
divisibility.
Thus both fixed lattices have rank one.
Choose primitive generators
$u_\psi$ and $u_{\psi^{-1}}$.
They are primitive in the full lattices:
each fixed kernel is saturated, since
the image of $S-I$ is torsion-free.
Semisimplicity over $K_3$ and preservation
of the pairing make the pairing of the
two fixed lines nonzero.

Unimodularity supplies
$w\in V_\psi$ with
$\langle w,u_{\psi^{-1}}\rangle=1$.
Writing the actual norm as
$\mathcal N=I+S+S^2$, we have
$\mathcal Nw=b\,u_\psi$ for some
$b\in\mathcal O$, and
\[
       b\,\langle u_\psi,u_{\psi^{-1}}\rangle
          =\langle\mathcal Nw,u_{\psi^{-1}}\rangle
          =3 .
                                                               \tag{D.71v}
\]
Fixed-cup divisibility says the
displayed pairing lies in $3\mathcal O$.
Since $\mathcal O$ is unramified,
(D.71v) forces that pairing to be
three times a unit and forces $b$
to be a unit.
The norm therefore maps $V_\psi$
onto its primitive fixed line.

Modulo three, the norm is
$(S-I)^2$, and its rank on this
three-dimensional block is one.
As $(S-I)^3=0$, the block is exactly
one Jordan block of size three.
This also supplies an actual integral
cyclic basis: lift a vector whose
three cyclic images form a basis
modulo three. The columns
$w,Sw,S^2w$ then have unit determinant,
so the map
$\mathcal O[C_3]\to V_\psi$ sending
the group generator to $S$ and $1$
to $w$ is an isomorphism.

For descent, no unspecified module
form is used.
After the splitting residue-field
extension, the norm on $W_3/3W_3$
has rank $q-1$ in dimension $3(q-1)$.
Ranks are unchanged by this field
extension. Since every Jordan block
has size at most three and each
size-three block contributes one
to the norm rank, all blocks over
$\F_3$ are of size three.
Choose a basis of the actual
$\F_3[C_3]$-module by cyclic
vectors and lift them to $W_3$.
The resulting explicit action map
\[
 \Z_3[C_3]^{q-1}\longrightarrow W_3,\qquad
             (f_i(S))_i\longmapsto\sum_i f_i(S)w_i
                                                               \tag{D.71w}
\]
reduces to an isomorphism modulo three.
It is a map between free $\Z_3$-modules
of the same finite rank, so its matrix
has unit determinant and is invertible.
Its inverse is $C_3$-equivariant.
This proves the integral freeness
in the statement.

\emph{The three full section classes and the fixed-kernel map.}
Put $E_N=N\rtimes\langle s\rangle$.
D.48 gives $H^1(C,\F_3)=0$,
so $\chi(C)=0$ and $\rho(C)\subset R$.
Its image is either $1$ or $R$.
The other global section class has
nonzero $\chi$-value by (D.71e),
so none of its conjugates meets $E_N$.
Every conjugate of $\langle s\rangle$
does meet $E_N$, because its image
in (D.71o) is the same central
order-three subgroup.
There is the exact class bijection
\[
 Q/\rho(C)\longrightarrow
 \{\text{$N$-conjugacy classes of sections in }E_N\},
 \qquad
 \rho(f)\rho(C)\longmapsto[f\langle s\rangle f^{-1}].
                                                               \tag{D.71x}
\]
Two conjugators give the same class
precisely when their ratio differs by
an element of $N$ and an element of
$N_F(\langle s\rangle)=C$.
Conjugacy under $E_N$ gives the same
classes, by removing the matching
section label from a conjugator.

Let $V=H^1(N,\F_3)$.
On $eV$ the action is two Jordan
blocks of size two, while (D.71w)
makes the complementary module free
over $\F_3[C_3]$.
The periodic cyclic resolution has
maps $S-I$ and
$I+S+S^2=(S-I)^2$.
Each size-two block contributes one
dimension in every positive degree;
the free blocks contribute zero.
Thus
\[
          \dim H^j(C_3,V)=2\quad(j>0),\qquad
          \dim H^n(E_N,\F_3)\leq3\quad(n\geq2).
                                                               \tag{D.71y}
\]
The last bound is the actual two-row
Hochschild--Serre sequence, since $N$
has cohomological dimension one.

Brown restriction applies to $E_N$:
$N$ is open of dimension one, and
every finite subgroup injects into
$C_3$, so there is no elementary
abelian subgroup of rank two.
The finite-dimensional source in
(D.71y) makes the dense restriction
embedding onto its Hausdorff target.
If $\rho(C)=1$, (D.71x) supplies
$3q$ normalizer factors, already
contradicting (D.71y).
Therefore $\rho(C)=R$, proving the
surjectivity in (D.71b).
The classes are exactly its three
cosets, represented by $1,t,t^2$.
Intersecting the full direct-product
normalizers with $E_N$ gives
exactly (D.71c).
The maps
$(h,s^j)\mapsto x^ihs^jx^{-i}$
and their inverses, which read
the fixed-group coordinate and the
specified quotient label after
conjugation by $x^{-i}$, prove
these identifications as full
topological groups.

Each fixed kernel $x^iC_0x^{-i}$
is a closed subgroup of the free
profinite group $N$ and has
$3$-cohomological dimension at most
one. Kunneth and Brown therefore give,
in high degrees,
\[
 3\bigl(1+\dim H^1(C_0,\F_3)\bigr)
                  =\dim H^n(E_N,\F_3)\leq3.
                                                               \tag{D.71z}
\]
All dimensions are nonnegative,
so $H^1(C_0,\F_3)=0$ and the
high-degree dimension is exactly
three. This completes the proof
of the fixed-kernel and section
assertions.

The executable certificate
\path{code/part_c/torus_cubic_frobenius_d71.py}
checks all eight based character
adaptations and their word inverses,
the integral covering matrices,
the Frobenius quotients and cusp
cycles for the first admissible
primes, the graph-cycle maps with
their inverses, and every rank-three
order-three matrix over $\F_3$.
The integral fixed-cup argument
and the full Brown normalizers
above supply the conclusions
beyond those finite identities.
\end{proof}
